Ce rapport présente la conception et la réalisation d'\textbf{IceVentes}, une application web de gestion des ventes de blocs de glace développée au cours d'un stage effectué au sein de l'Office National des Pêches (ONP). Dans le secteur de la pêche, la glace en blocs est indispensable à la conservation du poisson ; sa vente aux professionnels était jusqu'alors suivie de manière manuelle, sans centralisation ni traçabilité. Le projet vise à informatiser ce processus au moyen d'une application web moderne, sécurisée et simple d'utilisation.

La solution repose sur le cadriciel Python \textbf{Django}, une interface responsive construite avec Bootstrap~5 et HTMX, une base de données relationnelle (SQLite en développement, PostgreSQL en production) et des exports bureautiques (Excel, PDF). Elle couvre l'ensemble du cycle métier : authentification et gestion des utilisateurs avec trois rôles (Administrateur, Agent, Caissier) ; paramétrage d'un référentiel d'usines à glace (code, ville, prix unitaire en DH/Kg) ; gestion des demandeurs (clients, personnes physiques ou morales, acheteurs ou vendeurs) ; saisie d'une vente avec génération automatique d'un code unique \texttt{BG\_<réf>/<année> <numéro>} et calcul automatique de la quantité ou du prix total ; enfin, un historique multicritère avec exports et un tableau de bord synthétisant l'activité (nombre de ventes, quantité vendue, chiffre d'affaires, demandeurs actifs) ainsi qu'un graphique d'évolution mensuelle.

Une attention particulière a été portée à la cohérence des données (unicité des codes, référentiel actif unique, numérotation fiable en concurrence), à la sécurité des accès (contrôle par rôle, protection contre les attaques par force brute) et à la qualité de l'expérience utilisateur. Le rapport détaille le cadre du projet, la spécification des besoins, la conception du système et sa réalisation effective, illustrée par les interfaces produites.
